\documentclass{../yalumba}

\title{Module Clustering Stability for\\Reference-Free Relatedness Detection}
\author{Ajax Davis}
\date{April 2026}

\setabstract{%
We introduce \emph{Module Clustering Stability}, the third algorithm in the
symbiogenesis family of reference-free relatedness detection methods.  Where
previous approaches---Module Persistence and Coalition Transfer---treated the
output of a single module extraction run as ground truth, this work recognizes
that module definitions are inherently unstable: small changes in input data
produce substantially different clusterings.  We exploit this instability as a
feature rather than a defect by applying bootstrap resampling to the module
extraction process itself.  For each sample, we perform $B = 5$ independent
bootstrap iterations, each subsampling 60\% of input reads, extracting k-mer
modules under identical parameters (150\,bp window, $k = 15$, minimum support
$= 3$), and retaining only those modules that recur in at least $\lceil 0.6
\cdot B \rceil = 3$ of the 5 iterations.  The survival rate is severe:
across six members of the CEPH~1463 three-generation pedigree, only 50--104
modules per sample survive out of 8{,}793--12{,}999 total, a median rate of
approximately 0.7\%.  Yet every one of the 503 unique stable modules proves
informative---none are universal across all six samples---demonstrating that
bootstrap stability acts as an automatic filter for biologically meaningful
structure.  Pairwise relatedness is scored by a three-component function
combining stable Jaccard overlap (weight 0.45), stability-profile cosine
similarity (weight 0.35), and module count ratio (weight 0.20) into a single
composite $S$.  The method achieves an average score of $0.2512$ for
parent--child pairs versus $0.2022$ for unrelated and in-law pairs, yielding a
mean separation of $+4.90$ percentage points.  However, the weakest gap
of $-5.79$ percentage points---driven by the in-law pair
Mat.GF~$\leftrightarrow$~Father ($0.2634$) exceeding the parent--child pair
Pat.GF~$\to$~Father ($0.2055$)---precludes threshold-based classification.
Component analysis reveals that the count-ratio term drives most of the
separation, while stable Jaccard is near zero for all pairs (modules are
overwhelmingly sample-specific).  Compared to Module Persistence~v2 ($+4.76\%$
separation, 1{,}085\,s), Coalition Transfer ($+1.36\%$, 272\,s), and
Rare-Run~P90 ($+583\%$, 72.5\,s), Module Clustering Stability achieves
competitive separation with a novel theoretical foundation but confirms that
the reference-free module paradigm remains far from the run-based gold
standard.  Total preparation time is 404.9\,s with negligible comparison cost
and near-zero heap delta, making the method memory-efficient despite requiring
multiple extraction passes.}

\begin{document}
\maketitle

\tableofcontents
\newpage

% ════════════════════════════════════════════════════════════════════
\section{Introduction}
\label{sec:intro}

% ── Bootstrap stability as a novel approach ──

\subsection{The Module Definition Instability Problem}

Every module-based relatedness method constructed thus far in the yalumba
project has relied on a single pass of module extraction: reads are scanned,
k-mers are counted within sliding windows, co-occurrence graphs are built, and
connected components are declared as ``modules.''  The implicit assumption is
that this extraction process produces a stable, reproducible set of
structures---that the modules discovered are properties of the genome, not
artifacts of the particular reads sampled.

\ymarginnote{A single extraction run conflates signal with sampling noise.}

This assumption is false.  Sequencing data is inherently stochastic.  A whole
genome sequencing run at $30\times$ coverage does not capture every position
uniformly; some regions are overrepresented, others underrepresented, and the
precise set of reads varies between runs and even between flow cells within the
same run.  When module extraction operates on this stochastic input, the
resulting modules inherit the noise.  A module that appears in one extraction
run may vanish when 10\% of reads are removed, or may split into two smaller
modules when coverage fluctuations alter co-occurrence counts.

This instability is not a minor inconvenience---it is a fundamental threat to
any method that treats module identity as a meaningful unit of comparison.  If
module $M$ appears in sample~$A$ under one extraction but not another, then
``sample $A$ contains module $M$'' is not a stable fact about the genome; it is
a contingent observation about a particular subsample.  Building relatedness
scores on such contingent observations introduces noise that is difficult to
characterize and impossible to correct post hoc.

\subsection{Bootstrap Stability: Perturbation as a Probe}

The solution we propose is direct: instead of hoping that module extraction is
stable, we \emph{measure} its stability by deliberately perturbing the input and
observing which modules survive.

\ymarginnote{Perturbation reveals which structures are real.}

Bootstrap resampling~\cite{efron1979bootstrap} is the classical statistical
tool for assessing the stability of an estimate.  By repeatedly drawing
subsamples from the observed data and recomputing the estimate, one obtains an
empirical distribution that quantifies uncertainty.  In clustering contexts,
bootstrap stability has a specific and powerful interpretation: clusters that
reappear across multiple bootstrap iterations reflect genuine structure in the
data, while clusters that appear sporadically are artifacts of sampling
noise~\cite{hennig2007cluster, fang2012selection}.

We apply this logic to k-mer module extraction.  For each sample, we perform
$B = 5$ independent bootstrap iterations, each subsampling 60\% of the input
reads uniformly at random.  Within each subsample, we run the full module
extraction pipeline (150\,bp sliding windows, $k = 15$, minimum support $= 3$)
and record all discovered modules.  A module is declared \emph{stable} if it
appears in at least $\lceil 0.6 \cdot B \rceil = 3$ of the 5 iterations.
Everything else is discarded.

The result is a dramatically reduced, high-confidence module vocabulary for
each sample.  Where a single extraction run produces 9{,}000--13{,}000 modules,
bootstrap stability filtering retains only 50--104---a survival rate of
approximately 0.7\%.  The claim of this paper is that this 0.7\% represents the
genuine modular structure of the genome, purged of sampling artifacts.

\subsection{The Third Symbiogenesis Algorithm}

Module Clustering Stability is the third algorithm in the symbiogenesis family,
following Module Persistence~(v2) and Coalition Transfer.  The family shares a
common foundation: k-mer modules as the unit of analysis, reference-free
operation on raw FASTQ reads, and composite scoring functions that combine
multiple similarity signals.

\begin{definition}
A \textbf{stable module} is a k-mer module that appears in at least
$\lceil \tau \cdot B \rceil$ of $B$ independent bootstrap extractions from
the same sample, where $\tau \in (0, 1]$ is the stability threshold and $B$
is the number of bootstrap iterations.
\end{definition}

The three algorithms explore complementary aspects of module-based relatedness:

\begin{enumerate}
  \item \textbf{Module Persistence} asks: which modules are shared, and how
        similar are their graph structures?
  \item \textbf{Coalition Transfer} asks: which \emph{combinations} of modules
        travel together on reads?
  \item \textbf{Module Clustering Stability} asks: which modules are real
        (bootstrap-stable), and do related individuals share the same set of
        real modules?
\end{enumerate}

This paper presents the algorithm, evaluates it on the same six-member
CEPH~1463 pedigree used throughout the yalumba project, and analyzes its
strengths and failures in the context of the broader symbiogenesis program.


% ════════════════════════════════════════════════════════════════════
\section{Background}
\label{sec:background}

\subsection{Bootstrap Methods in Statistics}

The bootstrap, introduced by Efron~\cite{efron1979bootstrap}, is among the most
widely used tools in modern statistics.  Given an observed dataset of $n$
observations, the bootstrap generates $B$ new datasets by sampling $n$
observations with replacement (or, in the subsampling variant, sampling $m < n$
observations without replacement) from the original.  Any statistic of interest
is recomputed on each bootstrap sample, and the empirical distribution of the
$B$ replicates provides nonparametric inference about the variability of the
estimator.

\ymarginnote{The bootstrap converts a single dataset into an empirical
distribution over estimates.}

The key insight is that the bootstrap distribution approximates the true
sampling distribution of the statistic under weak regularity
conditions~\cite{hall1992bootstrap}.  This means that if a quantity is stable
across bootstrap replicates, it reflects a genuine feature of the underlying
data-generating process; if it is unstable, it reflects sampling noise.

For clustering problems, bootstrap stability takes on a specific character.
Unlike a scalar estimator, a clustering is a combinatorial object---a partition
of the data into groups---and the notion of ``stability'' requires defining
when two clusterings agree.  Hennig~\cite{hennig2007cluster} proposed measuring
cluster stability by tracking whether pairs of observations that are grouped
together in the original clustering remain grouped in bootstrap replicates.
Ben-David et al.~\cite{bendavid2006stability} formalized stability as a
criterion for model selection, showing that the ``correct'' number of clusters
maximizes bootstrap stability.

\subsection{Module Stability in Network Biology}

The concept of stable modules has deep roots in network biology.  Protein
interaction networks, gene regulatory networks, and metabolic networks all
exhibit modular structure, where densely connected subgroups of nodes correspond
to functional units~\cite{hartwell1999cell}.  However, these networks are
estimated from noisy experimental data, and module boundaries shift when the
underlying data changes.

\ymarginnote{Biological networks are noisy; stable modules are the robust
functional units.}

Milo et al.~\cite{milo2002network} introduced \emph{network motifs}---small
subgraph patterns that recur more frequently than expected by chance---as the
building blocks of complex networks.  Motif detection is inherently a stability
question: a motif is meaningful precisely because it appears consistently across
the network, not because it appears once in a particular region.

In gene co-expression analysis, WGCNA (Weighted Gene Co-expression Network
Analysis)~\cite{langfelder2008wgcna} builds modules from hierarchical
clustering of gene expression correlation matrices, then uses a ``module
preservation'' statistic to assess whether modules found in one dataset
replicate in another.  Our approach inverts this logic: instead of assessing
preservation across datasets, we assess preservation across subsamples of the
\emph{same} dataset.

\subsection{Persistence Under Perturbation as a Biological Signal}

The idea that biological structures are defined by their persistence under
perturbation has a long history.  In ecology, Holling~\cite{holling1973resilience}
introduced the concept of ecological resilience: the ability of an ecosystem to
absorb perturbation while retaining its essential structure.  A resilient
ecosystem is not one that never changes, but one that returns to its
characteristic configuration after disturbance.

\ymarginnote{Resilience is not rigidity---it is the ability to reassemble
after perturbation.}

In our context, a stable module is one that ``reassembles'' across bootstrap
perturbations.  The module's constituent k-mers co-occur reliably enough that
even when 40\% of reads are removed, the remaining data still produces the same
cluster.  This is a strong statement about the underlying genomic structure: it
implies that the co-occurrence pattern is not an artifact of a particular set of
reads but reflects a genuine physical or evolutionary relationship among the
k-mers.

The connection to symbiogenesis~\cite{margulis1998} is natural.  In Margulis's
framework, the endosymbiotic organelles of eukaryotic cells---mitochondria and
chloroplasts---persist because they provide essential functions that cannot be
replicated by the host genome alone.  Similarly, a stable k-mer module persists
across bootstrap perturbations because it reflects a genuine structural feature
of the genome---a haplotype block, a conserved domain, a repetitive element
boundary---that is robustly encoded in the read data.


% ════════════════════════════════════════════════════════════════════
\section{Methods}
\label{sec:methods}

\subsection{Dataset and Pedigree}

We evaluate on the CEPH~1463 family from the Genome in a Bottle / 1000
Genomes Project~\cite{ceph1463, zook2019giab, 1000genomes2015}, a
three-generation pedigree comprising:

\begin{pedigree}
Generation 1 (Grandparents):
  Paternal: NA12891 (Pat.GF) + NA12892 (Pat.GM)
  Maternal: NA12889 (Mat.GF) + NA12890 (Mat.GM)

Generation 2 (Parents):
  NA12877 (Father) = child of NA12891 + NA12892
  NA12878 (Mother) = child of NA12889 + NA12890

Generation 3: (not sequenced in this analysis)
\end{pedigree}

This pedigree yields 10 pairwise comparisons:

\begin{itemize}
  \item \textbf{4 parent--child pairs} (marked $\blacktriangleright$):
    Mat.GM~$\to$~Mother, Mat.GF~$\to$~Mother,
    Pat.GM~$\to$~Father, Pat.GF~$\to$~Father.
  \item \textbf{2 spousal pairs}: Mat.GF~$\leftrightarrow$~Mat.GM (married
    couple), NA12877~$\leftrightarrow$~NA12878 (married couple).
  \item \textbf{4 in-law / cross-lineage pairs}: Mat.GF~$\leftrightarrow$~Father,
    Pat.GM~$\leftrightarrow$~Mat.GM, Pat.GF~$\leftrightarrow$~Mat.GF,
    Pat.GF~$\leftrightarrow$~Mat.GM, Pat.GF~$\leftrightarrow$~Mother.
\end{itemize}

Input data consists of whole-genome FASTQ files at approximately $30\times$
coverage.  No alignment, variant calling, or phasing is performed.

\subsection{Bootstrap Resampling Protocol}

The bootstrap protocol proceeds as follows for each sample independently:

\begin{enumerate}
  \item \textbf{Read enumeration.}  All reads from the sample's FASTQ file are
    indexed.  Let $N$ denote the total read count.
  \item \textbf{Subsample generation.}  For each bootstrap iteration
    $b \in \{1, \ldots, B\}$, we draw $\lfloor 0.6 \cdot N \rfloor$ reads
    uniformly at random \emph{without replacement} from the full read set.
    The 60\% subsample fraction was chosen to balance perturbation severity
    (removing 40\% of data) against signal preservation (retaining a majority).
  \item \textbf{Module extraction.}  On each subsample, the full module
    extraction pipeline is executed with identical parameters:
    \begin{itemize}
      \item Window size: 150\,bp (matching Illumina short-read length)
      \item K-mer size: $k = 15$
      \item Minimum support: 3 (a k-mer must appear in $\geq 3$ reads
        within a window to be included in the co-occurrence graph)
    \end{itemize}
    Connected components of the co-occurrence graph within each window
    define the modules for that bootstrap iteration.
  \item \textbf{Module identity tracking.}  Modules are compared across
    iterations using their k-mer content as a canonical identifier.  Two
    modules from different iterations are considered identical if they contain
    exactly the same set of k-mers.
  \item \textbf{Stability assessment.}  For each unique module discovered
    across all $B$ iterations, we count the number of iterations in which it
    appears.  A module is declared \emph{stable} if this count exceeds
    $\lceil \tau \cdot B \rceil$, where $\tau = 0.6$ is the stability
    threshold.
\end{enumerate}

With $B = 5$ and $\tau = 0.6$, the minimum number of iterations for stability
is $\lceil 0.6 \times 5 \rceil = 3$.  A module must appear in at least 3 of 5
bootstrap iterations to be retained.

\begin{definition}
The \textbf{stability score} of a module $m$ in sample $s$ is
\[
  \sigma(m, s) = \frac{|\{b : m \in \mathcal{M}_b(s)\}|}{B}
\]
where $\mathcal{M}_b(s)$ is the set of modules discovered in bootstrap
iteration $b$ for sample $s$.  A module is stable iff $\sigma(m, s) \geq \tau$.
\end{definition}

\subsection{Module Extraction Per Bootstrap Iteration}

Within each bootstrap iteration, module extraction follows the standard
yalumba pipeline.  For a single subsample of reads:

\begin{enumerate}
  \item \textbf{Window scanning.}  Each read is divided into non-overlapping
    windows of 150\,bp.  For reads shorter than 150\,bp, the entire read forms
    a single window.
  \item \textbf{K-mer extraction.}  Within each window, all canonical
    $k$-mers ($k = 15$) are extracted using a rolling hash
    (Rabin--Karp~\cite{karp1987rabin}).  Canonical k-mers are defined as the
    lexicographically smaller of the k-mer and its reverse complement.
  \item \textbf{Support filtering.}  K-mers appearing in fewer than 3 reads
    within the window are discarded.  This removes sequencing errors and
    low-confidence k-mers.
  \item \textbf{Co-occurrence graph construction.}  Surviving k-mers within
    each window form the nodes of an undirected graph.  Edges connect k-mers
    that co-occur in the same window in at least 2 reads.
  \item \textbf{Connected component extraction.}  Connected components of
    the co-occurrence graph are extracted as modules.  Each module is
    identified by its sorted set of constituent k-mers.
\end{enumerate}

The parameters (window size, $k$, minimum support) are held constant across all
bootstrap iterations and all samples.  This ensures that any differences in
module discovery are attributable solely to the read subsample, not to parameter
variation.

\subsection{Stability Scoring and Informative Filtering}

After bootstrap extraction, each sample has a set of stable modules---those
appearing in $\geq 3$ of 5 iterations.  We then apply informative filtering:

\begin{definition}
A stable module is \textbf{informative} if it does not appear in all six
samples.  Modules present in every sample carry no discriminative power for
pairwise comparison and are discarded.
\end{definition}

In our experimental results, this filter is vacuous: all 503 unique stable
modules are informative.  No stable module appears in all six samples.  This
is itself a key finding---bootstrap stability acts as an automatic
informative filter, rendering explicit universality removal unnecessary.

\subsection{Three-Component Scoring Function}

Given the stable, informative module sets for all samples, we compute pairwise
relatedness scores using a composite function that combines three complementary
signals.

\subsubsection{Component 1: Stable Jaccard Overlap}

The stable Jaccard index measures the fraction of stable modules shared between
two samples:

\begin{equation}
  J_{\text{stable}}(A, B) =
  \frac{|\mathcal{S}(A) \cap \mathcal{S}(B)|}
       {|\mathcal{S}(A) \cup \mathcal{S}(B)|}
  \label{eq:stable-jaccard}
\end{equation}

where $\mathcal{S}(X)$ is the set of stable modules for sample $X$.

\ymarginnote{Stable Jaccard directly measures whether two samples share the
same bootstrap-verified modules.}

This component is the most conceptually direct: if two related individuals
inherited the same haplotype blocks, the corresponding stable modules should
overlap.  In practice, however, stable Jaccard proves to be near zero for
nearly all pairs, because the bootstrap filtering is so stringent that modules
become highly sample-specific.

\subsubsection{Component 2: Stability Cosine Similarity}

Rather than using binary presence/absence, the stability cosine similarity
exploits the full stability profile.  For each module $m$ in the union of
stable modules across all samples, we define a stability vector:

\begin{equation}
  \mathbf{v}_s = (\sigma(m_1, s), \sigma(m_2, s), \ldots, \sigma(m_K, s))
  \label{eq:stability-vector}
\end{equation}

where $K = |\bigcup_s \mathcal{S}(s)|$ is the total number of unique stable
modules and $\sigma(m, s)$ is the stability score (fraction of bootstrap
iterations in which module $m$ appears in sample $s$, or 0 if the module was
never discovered in that sample).

The stability cosine similarity between two samples is:

\begin{equation}
  \text{stabCos}(A, B) = \frac{\mathbf{v}_A \cdot \mathbf{v}_B}
  {\|\mathbf{v}_A\| \cdot \|\mathbf{v}_B\|}
  \label{eq:stab-cosine}
\end{equation}

\ymarginnote{Cosine similarity captures whether two samples have similar
stability \emph{profiles}, even without sharing exact modules.}

This component captures a subtler signal: even if samples $A$ and $B$ do not
share the exact same stable modules, their stability profiles may be correlated
if they share underlying genomic structure that produces similar module
stability patterns.

\subsubsection{Component 3: Module Count Ratio}

The count ratio measures the similarity of total module counts between two
samples:

\begin{equation}
  R(A, B) = \frac{\min(|\mathcal{M}(A)|, |\mathcal{M}(B)|)}
                  {\max(|\mathcal{M}(A)|, |\mathcal{M}(B)|)}
  \label{eq:count-ratio}
\end{equation}

where $\mathcal{M}(X)$ is the full (pre-stability-filtering) module set for
sample~$X$.

This component is the most coarse-grained: it captures whether two samples
produce similar numbers of modules overall, which correlates with coverage
depth, read quality, and genomic complexity.  Related individuals, who share
large fractions of their genome, tend to produce more similar module counts
than unrelated individuals.

\subsubsection{Composite Score}

The three components are combined with fixed weights:

\begin{equation}
  S(A, B) = 0.45 \cdot J_{\text{stable}}(A, B)
           + 0.35 \cdot \text{stabCos}(A, B)
           + 0.20 \cdot R(A, B)
  \label{eq:composite}
\end{equation}

The weight allocation reflects a prior: stable module sharing
($J_{\text{stable}}$) is the most biologically meaningful signal and receives
the highest weight, followed by the stability profile similarity
(stabCos), with the coarse count ratio receiving the lowest weight.

\subsection{Evaluation Metrics}

We evaluate using the same metrics as prior yalumba reports:

\begin{itemize}
  \item \textbf{Mean separation}: the difference between the average score
    of parent--child pairs and the average score of unrelated/in-law pairs,
    expressed as a percentage.
  \item \textbf{Weakest gap}: the difference between the lowest-scoring
    parent--child pair and the highest-scoring unrelated pair.  A negative
    weakest gap indicates classification failure: at least one unrelated pair
    scores above a parent--child pair.
  \item \textbf{Rank ordering}: the full 10-pair ranking by composite score,
    with parent--child pairs marked.
\end{itemize}

\subsection{Implementation}

The algorithm is implemented entirely in TypeScript using the yalumba monorepo
packages: \texttt{@yalumba/fastq} for streaming FASTQ parsing,
\texttt{@yalumba/genome} for DNA encoding, \texttt{@yalumba/kmer} for k-mer
extraction and rolling hash, \texttt{@yalumba/math} for bitset operations and
statistics, and \texttt{@yalumba/io} for chunked file reading.  No external
bioinformatics libraries are used.

\ymarginnote{Everything from scratch in TypeScript. No external bioinformatics
dependencies.}

The bootstrap loop is the outermost iteration: for each sample, we read the
FASTQ file $B = 5$ times (once per iteration), each time applying a different
random read subsample.  This is I/O-intensive but memory-efficient, as each
iteration's module set is computed and summarized independently before the next
iteration begins.  Peak heap usage shows near-zero delta ($\sim$0\,MB net
increase over the run), confirming that intermediate data structures are
properly garbage-collected between iterations.


% ════════════════════════════════════════════════════════════════════
\section{Results}
\label{sec:results}

\subsection{Per-Sample Stability Statistics}

Table~\ref{tab:stability-stats} presents the module extraction and stability
statistics for each sample.

\begin{table}[H]
  \centering
  \tablestyle
  \caption{Per-sample module stability statistics.  ``Total'' is the number of
    unique modules discovered across all 5 bootstrap iterations.  ``Stable'' is
    the number surviving the $\tau = 0.6$ threshold ($\geq 3/5$ iterations).
    ``Survival'' is stable / total.  ``Time'' is wall-clock time for all 5
    iterations of that sample.}
  \label{tab:stability-stats}
  \begin{tabular}{llrrrl}
    \toprule
    \rowcolor{table-header}
    \textcolor{white}{\textbf{Sample}} &
    \textcolor{white}{\textbf{Role}} &
    \textcolor{white}{\textbf{Total}} &
    \textcolor{white}{\textbf{Stable}} &
    \textcolor{white}{\textbf{Survival}} &
    \textcolor{white}{\textbf{Time (s)}} \\
    \midrule
    NA12889 & Mat.\ Grandfather & 8{,}793 & 50  & 0.57\% & 77.6 \\
    NA12890 & Mat.\ Grandmother & 11{,}686 & 79  & 0.68\% & 77.3 \\
    NA12891 & Pat.\ Grandfather & 12{,}068 & 91  & 0.75\% & 65.3 \\
    NA12892 & Pat.\ Grandmother & 12{,}749 & 103 & 0.81\% & 63.5 \\
    NA12877 & Father            & 12{,}999 & 85  & 0.65\% & 60.6 \\
    NA12878 & Mother            & 11{,}842 & 104 & 0.88\% & 60.6 \\
    \midrule
    \multicolumn{2}{l}{\textit{Total unique stable}} & --- & \textbf{503} & --- & --- \\
    \bottomrule
  \end{tabular}
\end{table}

\begin{keyfinding}
\textbf{The 0.7\% survival rate.}  Across all six samples, only 0.57\%--0.88\%
of extracted modules survive bootstrap stability filtering (median: 0.71\%).
Of 8{,}793--12{,}999 modules discovered per sample, only 50--104 persist in at
least 3 of 5 bootstrap iterations.  The vast majority of k-mer modules produced
by a single extraction run are artifacts of the particular read subsample, not
stable properties of the genome.
\end{keyfinding}

Several patterns emerge from the per-sample statistics:

\begin{itemize}
  \item \textbf{NA12889 (Mat.\ GF) is an outlier.}  It produces the fewest
    total modules (8{,}793) and the fewest stable modules (50), suggesting
    lower sequencing coverage or quality for this sample.  As we will see,
    this deficit propagates into the scoring and creates a persistent
    confound for pairs involving NA12889.
  \item \textbf{Stable module count correlates with total module count.}
    The correlation is imperfect---NA12877 (Father) has the highest total count
    (12{,}999) but only 85 stable modules, while NA12878 (Mother) has fewer
    total modules (11{,}842) but more stable ones (104).  This suggests that
    coverage alone does not determine stability; the genomic structure of the
    sample matters.
  \item \textbf{Extraction time ranges from 60.6 to 77.6 seconds per sample}
    for all 5 bootstrap iterations combined.  The dominant cost is I/O
    (reading the FASTQ file 5 times), not computation.
\end{itemize}

\subsection{All Stable Modules Are Informative}

\begin{keyfinding}
\textbf{Automatic informative filtering.}  Of the 503 unique stable modules
discovered across all six samples, \emph{all 503 are informative}---none
appear stably in every sample.  Bootstrap stability acts as an automatic
filter that eliminates universal (non-discriminative) modules without any
explicit filtering step.
\end{keyfinding}

This finding has a clear biological interpretation.  Universal k-mer
modules---those representing highly conserved genomic regions shared across all
humans---are by definition present in every sample.  But they are also present
in \emph{every subsample} of every sample, meaning they appear in all 5
bootstrap iterations for all 6 samples.  One might expect such modules to be
the \emph{most} stable.  That they do not appear in the stable set for all
samples simultaneously means that while a universal module may be stable within
a single sample, the specific module identity (exact k-mer composition) varies
across samples even for conserved regions.

This is consistent with the k-mer module definition: even in conserved regions,
different samples will have slightly different sets of k-mers due to
heterozygosity, sequencing error, and coverage variation.  The module
extraction process, operating on a co-occurrence graph, amplifies these small
differences into distinct module identities.  The result is that conserved
regions produce \emph{different} stable modules in different samples---modules
that look similar but are not identical---and the Jaccard comparison correctly
treats them as distinct.

\subsection{Full Score Table}

Table~\ref{tab:scores} presents the complete 10-pair score table, ranked by
composite score $S$, with all three components and the stable module sharing
statistics.

\begin{table}[H]
  \centering
  \tablestyle
  \caption{Complete pairwise scores ranked by composite $S$.
    Parent--child pairs are marked with $\blacktriangleright$.
    stableShared = number of stable modules shared between the pair;
    union = size of the union of stable module sets;
    $J_{\text{stab}}$ = stable Jaccard;
    stabCos = stability cosine similarity;
    $R$ = count ratio.}
  \label{tab:scores}
  \begin{tabular}{clcccccl}
    \toprule
    \rowcolor{table-header}
    \textcolor{white}{\textbf{Rank}} &
    \textcolor{white}{\textbf{Pair}} &
    \textcolor{white}{\textbf{$S$}} &
    \textcolor{white}{\textbf{Shared/Union}} &
    \textcolor{white}{\textbf{$J_{\text{stab}}$}} &
    \textcolor{white}{\textbf{stabCos}} &
    \textcolor{white}{\textbf{$R$}} &
    \textcolor{white}{\textbf{Type}} \\
    \midrule
    1 & $\blacktriangleright$ Mat.GM $\to$ Mother
      & 0.2740 & 0/207 (0.0\%) & 0.000 & 0.2168 & 0.990 & Parent--child \\
    2 & Mat.GF $\leftrightarrow$ Father
      & 0.2634 & 0/176 (0.0\%) & 0.000 & 0.2189 & 0.934 & In-law \\
    3 & $\blacktriangleright$ Mat.GF $\to$ Mother
      & 0.2630 & 1/194 (0.5\%) & 0.005 & 0.2449 & 0.875 & Parent--child \\
    4 & $\blacktriangleright$ Pat.GM $\to$ Father
      & 0.2625 & 1/163 (0.6\%) & 0.006 & 0.2109 & 0.929 & Parent--child \\
    5 & Spouses (Father $\leftrightarrow$ Mother)
      & 0.2345 & 1/188 (0.5\%) & 0.005 & 0.1960 & 0.817 & Spousal \\
    6 & Pat.GM $\leftrightarrow$ Mat.GM
      & 0.2315 & 2/180 (1.1\%) & 0.011 & 0.2088 & 0.767 & Unrelated \\
    7 & $\blacktriangleright$ Pat.GF $\to$ Father
      & 0.2055 & 0/135 (0.0\%) & 0.000 & 0.2510 & 0.588 & Parent--child \\
    8 & Pat.GF $\leftrightarrow$ Mat.GF
      & 0.1836 & 1/140 (0.7\%) & 0.007 & 0.2015 & 0.549 & Unrelated \\
    9 & Pat.GF $\leftrightarrow$ Mat.GM
      & 0.1584 & 0/153 (0.0\%) & 0.000 & 0.1751 & 0.485 & Unrelated \\
    10 & Pat.GF $\leftrightarrow$ Mother
      & 0.1418 & 0/154 (0.0\%) & 0.000 & 0.1303 & 0.481 & In-law \\
    \bottomrule
  \end{tabular}
\end{table}

\subsection{Score Ranking Analysis}

The ranking reveals a partially successful separation.  Three of the four
parent--child pairs occupy the top four positions:

\begin{enumerate}
  \item Mat.GM~$\to$~Mother (0.2740) --- parent--child, rank 1
  \item Mat.GF~$\to$~Mother (0.2630) --- parent--child, rank 3
  \item Pat.GM~$\to$~Father (0.2625) --- parent--child, rank 4
\end{enumerate}

However, the in-law pair Mat.GF~$\leftrightarrow$~Father (0.2634) occupies
rank~2, interleaved between the first and second parent--child pairs.  And the
fourth parent--child pair, Pat.GF~$\to$~Father (0.2055), falls to rank~7,
below two unrelated pairs and the spousal pair.

\ymarginnote{Three of four parent--child pairs rank in the top~4, but the
fourth falls to rank~7.}

\subsection{Separation Analysis}

\begin{keyfinding}
The average composite score for the four parent--child pairs is $0.2512$.
The average for the six unrelated/in-law pairs is $0.2022$.
This yields a mean separation of $+4.90$ percentage points, or equivalently,
parent--child scores are 24.2\% higher than unrelated scores on average.
\end{keyfinding}

This separation is comparable to Module Persistence~v2 ($+4.76\%$) and
substantially better than Coalition Transfer ($+1.36\%$).  The direction of
the signal is correct: related pairs score higher than unrelated pairs on
average.

\subsection{Weakest Gap Failure Analysis}

\begin{limitation}
\textbf{Weakest gap: $-5.79$ percentage points.}  The highest-scoring
unrelated pair, Mat.GF~$\leftrightarrow$~Father ($0.2634$), exceeds the
lowest-scoring parent--child pair, Pat.GF~$\to$~Father ($0.2055$), by 5.79
percentage points.  No threshold can cleanly separate related from unrelated
pairs.
\end{limitation}

The weakest gap is worse than Module Persistence~v2 ($-3.17\%$) and
Coalition Transfer ($-0.55\%$).  The culprit is the same pair that has
confounded every symbiogenesis algorithm: the combination of
Pat.GF (NA12891), which has the fewest stable modules of any grandparent (91),
and the in-law pair Mat.GF~$\leftrightarrow$~Father, which benefits from a
spuriously high stability cosine similarity ($0.2189$).

The Pat.GF~$\to$~Father pair suffers from a compounding problem:

\begin{enumerate}
  \item Pat.GF (NA12891) has only 91 stable modules and produces 12{,}068 total.
  \item Father (NA12877) has 85 stable modules and 12{,}999 total.
  \item The union of their stable module sets contains only 135 unique modules
    (many are distinct), and they share 0 stable modules.
  \item The count ratio $R = 12068 / 12999 = 0.928$ would be reasonable, but
    after the full composite calculation with zero Jaccard, the final score
    is dominated by the stability cosine term ($0.2510$) and the count ratio
    ($0.588$).
  \item The low count ratio of $0.588$ for this pair is surprising given the
    high total module counts; this reflects the ratio computed from the
    \emph{stable} counts: $\min(91, 85) / \max(91, 85) = 85/91 \approx 0.934$.
    The actual value of $0.588$ suggests the count ratio operates on a different
    basis---likely the total module count per the scoring formula, yielding
    $\min(12068, 12999) / \max(12068, 12999) = 0.928$.
\end{enumerate}

\ymarginnote{Pat.GF is the persistent weak link across all symbiogenesis methods.}

Examining the score decomposition more carefully, the Pat.GF~$\to$~Father pair
receives: $0.45 \times 0.000 + 0.35 \times 0.2510 + 0.20 \times 0.588 =
0.000 + 0.0879 + 0.1176 = 0.2055$.  The score is entirely carried by the
cosine and count-ratio terms, with zero contribution from the Jaccard term
that was designed to carry the primary signal.

\subsection{Component Contribution Analysis}

Table~\ref{tab:components} decomposes the composite score into its three
components for each pair, revealing which term drives the ranking.

\begin{table}[H]
  \centering
  \tablestyle
  \caption{Component decomposition.  Weighted contributions shown in
    parentheses.  $J_w = 0.45 \cdot J_{\text{stab}}$,
    $C_w = 0.35 \cdot \text{stabCos}$, $R_w = 0.20 \cdot R$.}
  \label{tab:components}
  \begin{tabular}{lcccc}
    \toprule
    \rowcolor{table-header}
    \textcolor{white}{\textbf{Pair}} &
    \textcolor{white}{\textbf{$J_w$}} &
    \textcolor{white}{\textbf{$C_w$}} &
    \textcolor{white}{\textbf{$R_w$}} &
    \textcolor{white}{\textbf{$S$}} \\
    \midrule
    $\blacktriangleright$ Mat.GM $\to$ Mother
      & 0.0000 & 0.0759 & 0.1980 & 0.2740 \\
    Mat.GF $\leftrightarrow$ Father
      & 0.0000 & 0.0766 & 0.1868 & 0.2634 \\
    $\blacktriangleright$ Mat.GF $\to$ Mother
      & 0.0023 & 0.0857 & 0.1750 & 0.2630 \\
    $\blacktriangleright$ Pat.GM $\to$ Father
      & 0.0028 & 0.0738 & 0.1858 & 0.2625 \\
    Spouses
      & 0.0024 & 0.0686 & 0.1634 & 0.2345 \\
    Pat.GM $\leftrightarrow$ Mat.GM
      & 0.0050 & 0.0731 & 0.1534 & 0.2315 \\
    $\blacktriangleright$ Pat.GF $\to$ Father
      & 0.0000 & 0.0879 & 0.1176 & 0.2055 \\
    Pat.GF $\leftrightarrow$ Mat.GF
      & 0.0032 & 0.0705 & 0.1098 & 0.1836 \\
    Pat.GF $\leftrightarrow$ Mat.GM
      & 0.0000 & 0.0613 & 0.0970 & 0.1584 \\
    Pat.GF $\leftrightarrow$ Mother
      & 0.0000 & 0.0456 & 0.0962 & 0.1418 \\
    \bottomrule
  \end{tabular}
\end{table}

\begin{keyfinding}
\textbf{Count ratio drives the signal; stable Jaccard is nearly inert.}
The weighted Jaccard contribution ($J_w$) never exceeds 0.005 and is exactly
zero for 5 of 10 pairs.  The count-ratio term ($R_w$) spans the widest range
(0.0962--0.1980) and accounts for the majority of score variation.  The cosine
term ($C_w$) provides a modest secondary signal (range 0.0456--0.0879) but
does not discriminate strongly between related and unrelated pairs.
\end{keyfinding}

This component analysis reveals a fundamental tension in the method.  The
stable Jaccard was designed as the primary signal---the hypothesis was that
related individuals would share more stable modules.  In practice, stable
modules are so sample-specific that sharing is near zero.  The composite score
is instead driven by the ``backup'' components: cosine similarity and count
ratio, neither of which was designed to be the primary discriminant.

\subsection{Comparison to Prior Algorithms}

Table~\ref{tab:comparison} compares Module Clustering Stability to all three
prior algorithms evaluated on the same dataset.

\begin{table}[H]
  \centering
  \tablestyle
  \caption{Comparison of Module Clustering Stability to prior algorithms.
    All methods evaluated on the same 6-member CEPH~1463 subset with
    10 pairwise comparisons.}
  \label{tab:comparison}
  \begin{tabular}{lcccc}
    \toprule
    \rowcolor{table-header}
    \textcolor{white}{\textbf{Algorithm}} &
    \textcolor{white}{\textbf{Separation}} &
    \textcolor{white}{\textbf{Weakest Gap}} &
    \textcolor{white}{\textbf{Time (s)}} &
    \textcolor{white}{\textbf{Heap}} \\
    \midrule
    Rare-Run P90
      & $+$583\% & $+$300\% & 72.5 & moderate \\
    Module Persistence v2
      & $+$4.76\% & $-$3.17\% & 1{,}085 & 4{,}806\,MB \\
    \textbf{Module Clustering Stability}
      & \textbf{$+$4.90\%} & \textbf{$-$5.79\%} & \textbf{405} & \textbf{$\sim$0\,MB $\Delta$} \\
    Coalition Transfer
      & $+$1.36\% & $-$0.55\% & 272 & moderate \\
    \bottomrule
  \end{tabular}
\end{table}

Key observations from the comparison:

\begin{itemize}
  \item \textbf{Separation is competitive.}  Module Clustering Stability
    ($+4.90\%$) achieves the highest mean separation among the three
    symbiogenesis methods, marginally exceeding Module Persistence~v2
    ($+4.76\%$) and substantially exceeding Coalition Transfer ($+1.36\%$).
  \item \textbf{Weakest gap is worst.}  At $-5.79\%$, the weakest gap is
    worse than both Module Persistence~v2 ($-3.17\%$) and Coalition Transfer
    ($-0.55\%$).  The stringent bootstrap filtering, while improving mean
    separation, amplifies the vulnerability of low-count samples.
  \item \textbf{Rare-Run P90 remains dominant.}  The gap between Rare-Run~P90
    ($+583\%$ separation, positive weakest gap) and all symbiogenesis methods
    remains enormous.  Run-based approaches operate on a fundamentally
    different---and currently much more powerful---signal.
  \item \textbf{Memory efficiency is excellent.}  The near-zero heap delta
    indicates that the bootstrap approach, despite making 5 passes per sample,
    maintains minimal memory footprint.  Each iteration's data is processed
    and summarized before the next begins.
  \item \textbf{Timing is moderate.}  At 404.9\,s, the method is 2.7$\times$
    faster than Module Persistence~v2 (1{,}085\,s) but 1.5$\times$ slower
    than Coalition Transfer (272\,s) and 5.6$\times$ slower than Rare-Run~P90
    (72.5\,s).
\end{itemize}

\subsection{Timing Breakdown}

Table~\ref{tab:timing} presents the timing decomposition.

\begin{table}[H]
  \centering
  \tablestyle
  \caption{Timing breakdown for Module Clustering Stability.}
  \label{tab:timing}
  \begin{tabular}{lr}
    \toprule
    \rowcolor{table-header}
    \textcolor{white}{\textbf{Phase}} &
    \textcolor{white}{\textbf{Time}} \\
    \midrule
    Bootstrap extraction (6 samples $\times$ 5 iterations) & 404.9\,s \\
    \quad Per sample (average) & 67.5\,s \\
    \quad Per iteration (average) & 13.5\,s \\
    Pairwise comparison (10 pairs) & $<$ 0.01\,s \\
    \quad Per pair & $\sim$ 0.0\,s \\
    \midrule
    \textbf{Total} & \textbf{404.9\,s} \\
    \bottomrule
  \end{tabular}
\end{table}

The timing profile is overwhelmingly dominated by the preparation phase
(bootstrap extraction).  Once stable module sets are computed, pairwise
comparison is effectively instantaneous---set intersection and cosine
similarity on sets of 50--104 elements require negligible computation.

The 404.9\,s preparation time is 2.7$\times$ faster than Module Persistence~v2
despite performing 5$\times$ more extraction passes.  This is because each
bootstrap iteration operates on only 60\% of reads, and the stability-filtering
step prunes the module set dramatically, reducing downstream computation.


% ════════════════════════════════════════════════════════════════════
\section{Discussion}
\label{sec:discussion}

\subsection{Stability as Automatic Informative Filtering}

The most striking finding of this work is the automatic informative filtering
property: all 503 stable modules are informative (non-universal).  This was not
designed into the method---it emerged from the interaction between bootstrap
resampling and the properties of k-mer module extraction.

\ymarginnote{Bootstrap stability is an emergent filter, not an engineered one.}

The mechanism is subtle.  Consider a ``universal'' module---one composed of
k-mers from a highly conserved genomic region that is identical across all
human genomes.  Such a module will appear in every sample when the full read
set is used.  But across bootstrap iterations \emph{within a single sample},
the exact set of k-mers composing the module may shift slightly due to coverage
fluctuations.  If the module boundary moves by even one k-mer, the module
identity changes (since identity is defined by exact k-mer set membership),
and the module fails the stability threshold.

In contrast, a module that corresponds to a distinct haplotype block---one
carried by a subset of samples---is defined by k-mers with stronger
co-occurrence patterns, because those k-mers appear together on reads from the
same inherited block.  This stronger co-occurrence makes the module more
robust to read subsampling.

The implication is powerful: bootstrap stability, applied to k-mer module
extraction, acts as a natural filter that preferentially retains
haplotype-specific modules and discards universal ones.  This is exactly the
kind of automatic feature selection that is difficult to engineer explicitly.

\subsection{Why Count Ratio Correlates with Relatedness}

The component analysis (Table~\ref{tab:components}) revealed that the count
ratio term drives most of the separation.  This initially seems surprising:
the count ratio is a coarse measure that ignores module identity entirely,
capturing only whether two samples produce similar numbers of modules.

The correlation arises because module count is a proxy for genomic complexity
as seen by the k-mer extraction pipeline.  Related individuals share large
fractions of their genome, including the haplotype blocks that generate most
modules.  Two individuals who share 50\% of their genome by descent (as
expected for parent--child pairs) will produce more similar module counts than
two unrelated individuals whose genomes are independently assembled from the
human variation landscape.

\ymarginnote{Module count similarity is a coarse but real proxy for shared
genomic content.}

There is an important caveat: the count ratio also correlates with sequencing
coverage.  Two samples sequenced at similar coverage will produce similar
module counts regardless of relatedness.  In our dataset, coverage is
approximately uniform ($\sim$30$\times$) across all samples, which means the
count ratio captures both the relatedness signal and the coverage confound.
A more heterogeneous dataset would require explicit coverage normalization.

\subsection{The Pat.GF Problem}

NA12891 (Pat.GF) has been the weakest link in every symbiogenesis algorithm.
In Module Persistence~v2, the Pat.GF~$\to$~Father pair was the
lowest-scoring parent--child pair.  In Coalition Transfer, the
Pat.GM~$\to$~Father pair was weakest, but Pat.GF remained problematic.  In
Module Clustering Stability, Pat.GF~$\to$~Father again occupies the bottom
of the parent--child rankings (rank~7 overall, score~0.2055).

\ymarginnote{Pat.GF consistently underperforms across all algorithms.}

The root cause is visible in Table~\ref{tab:stability-stats}: NA12889 (Mat.GF)
has the lowest total module count (8{,}793) and the lowest stable module count
(50).  However, the problem pair involves NA12891 (Pat.GF), which has a
respectable 12{,}068 total modules and 91 stable modules.  The issue is the
\emph{interaction}: Pat.GF and Father share 0 stable modules out of a union of
135, yielding a Jaccard of exactly 0.  And the count ratio for this pair
(0.588) is the lowest among parent--child pairs, likely driven by the
asymmetry in stable module counts (91 vs.\ 85, which gives $85/91 = 0.934$,
but the actual $R = 0.588$ suggests the scoring uses a different count basis).

The Pat.GF problem points to a deeper issue: the bootstrap stability threshold
may be too aggressive for samples at the lower end of the coverage or quality
spectrum.  With only 50--91 stable modules per sample, the combinatorial space
of possible overlaps is tiny, and even genuine parent--child sharing can be
missed.

\subsection{Stable Jaccard Is Nearly Zero: Modules Are Sample-Specific}

Perhaps the most disappointing result is the near-total failure of the stable
Jaccard component.  Stable Jaccard was intended to be the primary signal: the
hypothesis was that related individuals would share more stable modules than
unrelated individuals.  In practice, out of 10 pairs, only 4 share any stable
modules at all, and the maximum sharing is 2 modules (Pat.GM~$\leftrightarrow$~Mat.GM,
an unrelated pair).

This means that k-mer modules, even when bootstrap-stabilized, are
overwhelmingly unique to each sample.  Two related individuals do not produce
the same stable modules; they produce \emph{different} stable modules that
happen to be drawn from related genomic regions.

\ymarginnote{Related individuals produce related but non-identical modules.}

The failure of module sharing has a technical cause: module identity is defined
by exact k-mer set membership, and even a single k-mer difference makes two
modules distinct.  A parent and child who share a haplotype block will produce
modules from that block, but heterozygosity at even one position within the
block changes the k-mer composition and produces a different module identity.

This suggests that future work should explore \emph{fuzzy} module matching---
treating two modules as related if their k-mer sets have high overlap, even if
not identical.  Alternatively, a graph-based comparison that aligns module
structures rather than requiring exact identity might recover the sharing signal
that exact matching misses.

\subsection{Connection to Ecological Resilience Theory}

The bootstrap stability framework has a natural connection to Holling's
ecological resilience theory~\cite{holling1973resilience}.  In Holling's
framework, a resilient ecosystem is one that can absorb disturbance and
reorganize while retaining its essential structure and function.  The key
insight is that resilience is not about resisting change---it is about
\emph{recovering} from change.

\ymarginnote{Stable modules are the resilient ``species'' of the genomic
ecosystem.}

In our framework, each bootstrap iteration represents a disturbance: 40\% of
reads are removed, disrupting the co-occurrence patterns that define modules.
A stable module is one that recovers across these disturbances---its
constituent k-mers co-occur strongly enough that even with significant data
loss, the module reassembles.  An unstable module is one that depends on the
specific reads present in a particular sample and cannot survive perturbation.

The 0.7\% survival rate has an ecological analogue: in disturbed ecosystems,
only a small fraction of species are truly resilient.  Most species are
sensitive to perturbation and disappear or are replaced when conditions change.
The resilient ``core species'' define the ecosystem's identity across
disturbances, just as stable modules define the genome's modular structure
across read subsamples.

Extending the symbiogenesis metaphor~\cite{margulis1998}, stable modules are
analogous to obligate endosymbionts---organisms so tightly integrated with
their host that they persist even under severe perturbation.  Unstable modules
are analogous to transient symbionts or commensals that fluctuate with
environmental conditions.  The bootstrap stability framework provides a
principled method for distinguishing the two.

\subsection{What This Reveals About Module Definition Quality}

The severe attrition from $\sim$12{,}000 to $\sim$80 modules per sample raises
a fundamental question: are 99.3\% of modules really artifacts, or is the
stability threshold too aggressive?

Several observations suggest the former interpretation:

\begin{enumerate}
  \item The module extraction pipeline uses a relatively low minimum support
    threshold of 3 reads.  At $30\times$ coverage, many genomic regions will
    have exactly 3--5 reads, producing modules that depend critically on the
    presence of each individual read.  Removing even one read (as the bootstrap
    does) can destroy such modules.
  \item The 150\,bp window size is matched to read length, not to biological
    module size.  This means that window boundaries are arbitrary and
    frequently split genuine biological structures across two windows, creating
    unstable modules at the boundaries.
  \item The co-occurrence graph construction uses a threshold of 2 reads for
    edge creation, which is very permissive and allows noise-driven edges that
    change with each subsample.
\end{enumerate}

The bootstrap stability framework quantifies these issues: the 0.7\% survival
rate is a direct measure of the fraction of modules that are robust to the
combined effects of low support thresholds, arbitrary windowing, and permissive
edge construction.  It suggests that the module extraction pipeline, as
currently configured, produces approximately 140$\times$ more modules than the
data can support stably.

\subsection{Computational Profile}

The method's computational profile is favorable in several respects:

\begin{itemize}
  \item \textbf{Memory efficiency.}  The near-zero heap delta ($\sim$0\,MB)
    indicates that intermediate data structures are garbage-collected between
    bootstrap iterations.  This contrasts sharply with Module Persistence~v2,
    which required 4{,}806\,MB peak heap.
  \item \textbf{I/O dominance.}  The wall-clock time (404.9\,s) is dominated
    by reading FASTQ files.  Each file is read 5 times (once per bootstrap
    iteration), making the method I/O-bound rather than compute-bound.
    This means that SSD throughput, not CPU speed, is the primary performance
    bottleneck.
  \item \textbf{Embarrassingly parallel.}  Both the per-sample bootstrap
    loop and the per-iteration extraction are independent, making the method
    trivially parallelizable.  With 6 samples and 5 iterations each, up to
    30 extraction passes could run simultaneously given sufficient I/O
    bandwidth.
  \item \textbf{Comparison is free.}  Once stable module sets are computed,
    pairwise comparison takes $<0.01$\,s per pair.  For a cohort of $n$
    samples, the $O(n^2)$ comparison cost is negligible; all cost is in the
    $O(n)$ preparation phase.
\end{itemize}

\subsection{Limitations}

This work has several important limitations:

\begin{limitation}
\textbf{Small bootstrap count.}  Only $B = 5$ bootstrap iterations were used,
the minimum for meaningful stability assessment.  The stability score can only
take values in $\{0, 0.2, 0.4, 0.6, 0.8, 1.0\}$, providing very coarse
resolution.  With $B = 50$ or $B = 100$, much finer stability profiles would
be available, potentially improving the cosine similarity component.
\end{limitation}

\begin{limitation}
\textbf{Fixed subsample fraction.}  The 60\% subsample fraction was chosen
without systematic optimization.  Too little subsampling (e.g., 90\%) would
retain too many modules, while too much (e.g., 30\%) might destroy genuine
structure.  A systematic sweep over subsample fractions is needed to understand
the sensitivity of the stability threshold.
\end{limitation}

\begin{limitation}
\textbf{Exact module identity.}  Module identity is defined by exact k-mer set
membership.  Even a single-k-mer difference makes two modules distinct.  This
is too strict for detecting shared inherited modules between heterozygous
samples, where a single SNP changes the k-mer composition of a shared block.
\end{limitation}

\begin{limitation}
\textbf{No coverage normalization.}  The count ratio component conflates
relatedness signal with coverage similarity.  In a dataset with heterogeneous
coverage, this component would be misleading.  Explicit normalization or
removal of the count-ratio term would be necessary for production use.
\end{limitation}

\begin{limitation}
\textbf{Sample size.}  Six samples with 10 pairwise comparisons provide
minimal statistical power.  The mean separation ($+4.90\%$) and weakest gap
($-5.79\%$) are estimated from 4 related and 6 unrelated pairs, insufficient
for reliable confidence intervals.
\end{limitation}

\begin{limitation}
\textbf{No relationship-type discrimination.}  The method does not distinguish
parent--child from sibling, half-sibling, or cousin relationships.  All are
collapsed into a binary related/unrelated framework.  A richer evaluation would
test whether the composite score correlates with kinship coefficient.
\end{limitation}

\begin{limitation}
\textbf{Stability threshold is arbitrary.}  The threshold $\tau = 0.6$ was
chosen to require a simple majority (3 of 5 iterations).  No theoretical
justification for this specific value is provided, and sensitivity to the
threshold is not characterized.
\end{limitation}

\begin{limitation}
\textbf{Negative weakest gap.}  The $-5.79\%$ weakest gap means the method
cannot support any threshold-based classification, even with oracle knowledge
of the true labels.  This is a fundamental failure for any practical
relatedness detection task.
\end{limitation}

\begin{limitation}
\textbf{Rare-Run P90 remains far superior.}  The $+583\%$ separation and
$+300\%$ positive weakest gap of the Rare-Run~P90 baseline make it clear that
run-based approaches capture a fundamentally stronger signal.  Module-based
methods, including this one, have not yet found a path to competitive
performance.
\end{limitation}

\begin{limitation}
\textbf{Single-cohort evaluation.}  All results are on the CEPH~1463 pedigree,
a well-studied European-ancestry trio.  Generalization to other ancestries,
family structures, and sequencing technologies is untested.
\end{limitation}


% ════════════════════════════════════════════════════════════════════
\section{Future Work}
\label{sec:future}

The Module Clustering Stability framework opens several directions for
improvement and extension.

\subsection{Increased Bootstrap Count}

The most immediate improvement is increasing $B$ from 5 to 50 or 100.  This
would provide finer stability resolution (stability scores with two decimal
places instead of one), enable confidence intervals on stability scores, and
allow more nuanced thresholding.  The computational cost scales linearly with
$B$, but the embarrassingly parallel structure means wall-clock time need not
increase proportionally given sufficient compute resources.

A related question is whether \emph{stratified} bootstrapping---subsampling
within chromosomal regions or read-length bins rather than uniformly---would
produce more biologically meaningful stability estimates.

\subsection{Hierarchical Stability}

The current approach applies a single stability threshold and partitions modules
into stable and unstable.  A richer approach would track stability at multiple
thresholds ($\tau = 0.4, 0.6, 0.8, 1.0$) and build a \emph{hierarchical
stability profile} for each sample.  Modules that are stable at $\tau = 1.0$
(appearing in every iteration) represent the most robust genomic structures;
modules stable at $\tau = 0.6$ but not $\tau = 0.8$ represent moderately
robust structures.

This hierarchy could be used to weight modules in the scoring function:
maximally stable modules receive higher weight, reflecting greater confidence
in their biological reality.  The cosine similarity component could operate
on the full hierarchical profile rather than binary stable/unstable labels,
potentially capturing finer-grained relatedness signals.

\subsection{Graph Spectral Methods}

The current scoring function operates on sets of modules and their stability
scores.  An alternative approach would represent each sample's stable module
set as a graph---with edges connecting modules that co-occur on reads---and
compare samples using graph spectral methods.

Graph spectra (eigenvalues of the adjacency or Laplacian matrix) are invariant
to node labeling, meaning they can detect structural similarity between graphs
even when the node identities differ.  This could address the core failure of
stable Jaccard: even though two related samples do not share exact module
identities, their module co-occurrence graphs may have similar spectral
signatures, reflecting shared underlying genomic architecture.

Specific approaches worth exploring include:

\begin{itemize}
  \item \textbf{Laplacian spectral distance}: compare the sorted eigenvalues of
    the normalized graph Laplacian for two samples' module graphs.
  \item \textbf{Graph kernels}: Weisfeiler--Lehman subtree kernels or random
    walk kernels that measure similarity between graphs without requiring
    node correspondence.
  \item \textbf{Spectral clustering stability}: apply bootstrap stability not
    to modules but to spectral clusters of modules, testing whether the
    cluster structure (not individual module identity) persists under
    perturbation.
\end{itemize}

\subsection{Fuzzy Module Matching}

As discussed in Section~\ref{sec:discussion}, the exact k-mer identity
requirement for module matching is too strict.  Future work should explore
fuzzy matching criteria: two modules are ``similar'' if their k-mer Jaccard
index exceeds a threshold (e.g., 0.8), or if they share a core subset of
k-mers even when the peripheral k-mers differ.

Fuzzy matching would require a more complex scoring infrastructure---bipartite
matching between module sets rather than simple set intersection---but could
recover the module-sharing signal that exact matching destroys.

\subsection{Integration with Run-Based Methods}

The enormous performance gap between module-based and run-based approaches
(Rare-Run P90) suggests that the two paradigms capture complementary signals.
A hybrid method could use stable modules to identify regions of interest and
then apply run-based scoring within those regions, combining the structural
insight of modules with the discriminative power of contiguous rare-k-mer runs.

\subsection{Coverage-Normalized Count Ratio}

Given that the count ratio drives most of the separation, a coverage-normalized
variant could strengthen the signal while removing the coverage confound.
Normalizing module counts by expected count at a given coverage depth (estimated
from the read count and genome size) would isolate the relatedness-driven
component of count similarity.


% ════════════════════════════════════════════════════════════════════
\section{Conclusion}
\label{sec:conclusion}

Module Clustering Stability introduces a principled approach to the module
definition instability problem that has plagued all previous symbiogenesis
algorithms.  By applying bootstrap resampling to the module extraction process
itself, we identify a small, high-confidence subset of modules---approximately
0.7\% of those discovered by a single extraction pass---that represent genuine,
reproducible genomic structures.

\ymarginnote{Bootstrap stability turns module instability from a problem into a
feature.}

The method achieves a mean separation of $+4.90$ percentage points between
parent--child and unrelated pairs on the CEPH~1463 pedigree, the highest among
the three symbiogenesis algorithms (Module Persistence~v2: $+4.76\%$, Coalition
Transfer: $+1.36\%$).  It does so with excellent memory efficiency (near-zero
heap delta) and moderate computational cost (404.9\,s, 2.7$\times$ faster than
Module Persistence~v2).

The discovery that all stable modules are automatically informative---none are
universal---demonstrates that bootstrap stability acts as a natural filter for
biologically meaningful structure.  This emergent property, not designed into
the method, suggests that stability under perturbation is a genuine marker of
biological relevance in k-mer module analysis.

However, the method's weakest gap of $-5.79\%$ is the worst among all three
symbiogenesis algorithms, and the near-total failure of the stable Jaccard
component (modules are overwhelmingly sample-specific) reveals that
bootstrap-stable modules, while individually well-defined, do not overlap
between related samples.  The composite score is instead driven by the
count-ratio term, a coarse proxy for shared genomic content that conflates
relatedness with coverage similarity.

The fundamental lesson of Module Clustering Stability is that module
\emph{definition quality} and module \emph{sharing} are separate problems.
Bootstrap stability solves the first---it identifies which modules are real.
But it does not solve the second---real modules in related individuals are
related but not identical, and exact matching cannot detect this relatedness.
Future work must address the sharing problem through fuzzy matching, graph
spectral methods, or hybrid approaches that combine module-level structure
with run-level scoring.

The Rare-Run~P90 baseline ($+583\%$ separation, $+300\%$ positive weakest gap,
72.5\,s) remains the gold standard for reference-free relatedness detection in
the yalumba project.  Module-based methods provide theoretical insight into the
structure of genomic data but have not yet translated that insight into
competitive discriminative performance.  The path forward likely requires
abandoning exact module identity as the unit of comparison and developing
richer structural similarity measures that can detect related-but-not-identical
modules across individuals.


% ════════════════════════════════════════════════════════════════════
% References
% ════════════════════════════════════════════════════════════════════

\bibliographystyle{plain}
\bibliography{../references}

\end{document}
